\chapter{Giriş}
\label{ch:giris}

\section{Problem}

Nesne takibi görevlerinde Meta'nın \emph{Segment Anything Model 2}
(SAM~2)~\cite{ravi2024sam2} yüksek doğruluk verir, ancak sunucu sınıfı
bir GPU dışında pratik değildir; edge cihazlarda saniyede tek kareye
kadar düşebilmektedir. \emph{EdgeTAM}~\cite{zhou2025edgetam}, SAM~2'nin
memory mekanizmasını yeniden tasarlayarak bu maliyeti telefon sınıfı
donanımda gerçek zamanlıya (iPhone 15 Pro Max'te 16~FPS) indirir.

İki modeli ayıran şey bir doğruluk-hız takasıdır:
\cref{fig:sam2-edgetam-takas}, EdgeTAM'ın SAM~2'ye yakın doğruluğu çok
daha yüksek bir kare hızında verdiğini gösterir. Bu çalışmanın
başlangıç noktası da budur: doğruluğu zaten kabul edilebilir olan bir
mimariyi alıp hedef donanımda gerçekten o hızda çalıştırmak.

\begin{figure}[H]
  \centering
  \gorselyertutucu{EdgeTAM makalesinden~\cite{zhou2025edgetam} alınan
    doğruluk-hız takası grafiği: SAM~2 ve EdgeTAM varyantlarının
    kare hızına karşı segmentasyon doğruluğu}
  \caption{SAM~2 ve EdgeTAM arasındaki doğruluk-hız takası.}
  \label{fig:sam2-edgetam-takas}
\end{figure}

Bu çalışmanın konusu EdgeTAM'ın kendisini değiştirmek değil, onun
eğitilmiş kontrol noktasını (checkpoint) \textbf{NVIDIA Jetson AGX
Orin} üzerinde TensorRT ile gerçek zamanlı çalışacak hale getirmektir.
Hedef tek cümleyle: doğruluktan ödün vermeden kare başına \textbf{40~ms}
(en az 25~FPS) bütçesinin altına inmek.

Başlangıç durumu, bu tür bir hızlandırmanın tipik ilk adımıydı:
yalnızca image encoder TensorRT'e taşınmış, diğer üç modül PyTorch'ta
kalmıştı. İlk taşınan modülün image encoder olması tesadüf değildir,
çünkü TensorRT'in istediği her özelliğe sahiptir: tek bir sabit
girdi/çıktı şekli, çağrılar arasında saklanan hiçbir durum ve veriye
bağlı hiçbir dallanma. Bu haliyle hiçbir yeniden yazım gerektirmeden
tek seferde ONNX'e aktarılır. Diğer üç modülde bu özelliklerin neden
bulunmadığı ve bunun nasıl çözüldüğü \cref{ch:sistem-yontem}'in
konusudur.

\section{Başlangıç Varsayımı ve Neden Yanlış Çıktığı}

Yalnızca image encoder'ın hızlandırılmış olması, yaygın bir varsayımın
sonucudur: en çok parametreyi taşıyan modül, zamanın çoğunu da harcıyor
olmalıdır. SAM~2 üzerinde bu varsayımla çalışan önceki işler, daha
küçük ve verimli bir image encoder tasarlamaya yoğunlaşmıştır. Sonuç
şuna benzer: motoru güçlendirilmiş bir araba, ama trafikte beklediği
için yine yavaş. Buradaki trafik memory attention'dır; geçerli kareyi
hafızadaki tüm geçmiş karelerle karşılaştırma maliyeti, encoder ne
kadar hızlanırsa hızlansın yerinde durur.

Bu çalışma varsayımı devralmak yerine ölçmüştür
(\cref{sec:flop-dagilimi}): kare başına işlem yükünün (FLOP)
\textbf{\%61,9}'u memory attention modülünde, yalnızca \textbf{\%26,8}'i
image encoder'dadır. EdgeTAM makalesi de aynı noktayı vurgular; memory
attention blokları bir gecikme darboğazıdır ve yalnızca image encoder'ı
küçültmek beklenen hızlanmayı getirmez.

\textbf{Neden bu ölçüm doğru sorunun cevabıdır.} Ölçülen şey süre
değil, FLOP'tur: gerçek kontrol noktası üzerinde, gerçek 1024×1024
girdi çözünürlüğünde, her modülün bir kare için yaptığı
çarpma-toplama sayısı (\texttt{tools/analyze\_cost.py}, PyTorch'un
kendi FLOP sayacıyla, CPU üzerinde). Bu tercih bilinçlidir. O aşamada
image encoder zaten TensorRT'te, diğer üç modül PyTorch'ta çalışıyordu;
kronometreyle yapılacak bir ölçüm büyük ölçüde \emph{hangi modülün
hangi arka uçta olduğunu} yansıtırdı, hesabın gerçekte nerede
yoğunlaştığını değil. Yani zaten hızlandırılmış modülün hızlı çıkması
kaçınılmazdı ve bu, hangi modülü hızlandırmaya değeceği sorusunu
cevaplamazdı. FLOP ise donanımdan, saat hızından, TensorRT sürümünden
ve arka uç seçiminden bağımsızdır; mimarinin talep ettiği aritmetiği
verir, dolayısıyla herhangi bir arka uçtan istenebilecek işin üst
sınırıdır. Bu sayede soru, hızlandırma çalışmasına başlamadan
\emph{önce} cevaplanabilmiştir.

Bu bulgu, çalışmanın kapsamını tek bir modülden EdgeTAM'ın
\textbf{dört} ana hesaplama bloğunun tamamına genişletmiştir.

\section{Katkılar}

Bu rapor şu katkıları belgelemektedir:

\begin{itemize}
  \item EdgeTAM'ın dört ana modülünün (image encoder, memory attention,
        memory encoder, SAM head) TensorRT'e aktarılabilir hale
        getirilmesi, bunun için sabit şekilli memory bank, karmaşık
        sayı içermeyen RoPE ve CUDA Graph ile tekrar oynatmanın
        kullanılması (\cref{ch:sistem-yontem});
  \item modüllerin \emph{neden} dört ayrı motor olarak bırakıldığının
        ölçümle gerekçelendirilmesi, tek motora birleştirmenin
        kazandıracağı sürenin doğrudan ölçülüp değmeyeceğine karar
        verilmesi (\cref{sec:fuzyon-karari});
  \item fp16'nın bf16'ya karşı hem doğruluk hem TensorRT desteği
        açısından üstün olduğunun ölçülmesi (\cref{sec:hassasiyet});
  \item her sayının hangi araçtan geldiğini sabitleyen bir ölçüm
        metodolojisi (\cref{sec:olcum-yontemi}).
\end{itemize}
